\chapter{Grief}

\section*{Definition}
A natural, multifaceted emotional, cognitive, and somatic response to loss -- most commonly bereavement -- characterized by sadness, yearning, preoccupation with the deceased, and impaired functioning. Prolonged grief disorder (PGD) is diagnosed when intense yearning persists beyond 12 months (6 months in children) and significantly impairs daily function.

\section*{What Happens in Your Body}
Grief activates the same neurocircuitry as attachment and reward: the anterior cingulate cortex, insula, and nucleus accumbens modulate the experience of separation distress. Dopaminergic reward prediction error from the loss of the attachment figure drives craving and preoccupation. Chronic grief dysregulates the hypothalamic-pituitary-adrenal (HPA) axis, producing sustained cortisol elevation and immune suppression, which contributes to the well-documented increase in cardiovascular morbidity during the first 6 months post-loss.

\section*{Causes}
\begin{itemize}
  \item \textbf{Primary:} Death of a loved one (spouse, child, parent, sibling, close friend)
  \item \textbf{Non-death losses:} Divorce or relationship dissolution, loss of a pet, job loss, loss of physical function or independence, loss of a home or community, amputation, loss of future expectations (e.g., infertility)
  \item \textbf{Complicated/prolonged:} Sudden or traumatic death, death of a child, lack of social support, concurrent major stressors, history of depression or anxiety disorders, insecure attachment style
\end{itemize}

\section*{When to See a Doctor}
See your doctor if:
\begin{itemize}
  \item Intense yearning, numbness, or identity disruption persists beyond 12 months
  \item Suicidal ideation, intent, or plan
  \item Inability to perform basic self-care (eating, hygiene, medication adherence)
  \item Concurrent major depressive episode, panic attacks, or substance use disorder
  \item Somatic complaints (chest pain, palpitations) requiring medical clearance
\end{itemize}

\section*{Related Symptoms}
\begin{itemize}
  \item Major depressive disorder (depressed mood, anhedonia, worthlessness -- more pervasive than grief)
  \item Post-traumatic stress disorder (intrusions, avoidance, hyperarousal following traumatic death)
  \item Anxiety disorders (generalized anxiety, panic)
  \item Insomnia and appetite changes
  \item Substance use
\end{itemize}
\textbf{Important:} Grief is not a mental disorder per se; the key distinction from major depression is that grief features intense yearning and episodic waves of pain while preserving the capacity for positive emotions and pleasure.

\section*{Self-Care / Home Management}
\begin{itemize}
  \item Allow unstructured time for emotional expression; avoid suppression or avoidance
  \item Maintain daily routines for sleep, meals, and physical activity as anchors
  \item Seek connection: support groups, trusted friends, religious or spiritual communities
  \item Limit major life decisions for at least 3--6 months after the loss
  \item Journaling or legacy projects (photo albums, letters, memory boxes) can aid processing
\end{itemize}

\section*{How Your Doctor Will Evaluate You}
\begin{itemize}
  \item Clinical interview: nature of the loss, time course, functional impact, prior mental health history
  \item Screen for suicidal ideation and risk factors
  \item Use validated instruments (PG-13 for prolonged grief, PHQ-9 for depression, GAD-7 for anxiety)
  \item Rule out medical causes of somatic symptoms (chest pain, weight loss, fatigue -- basic labs)
  \item Assess social support system and coping resources
\end{itemize}

\section*{Treatment Options}
\begin{itemize}
  \item Supportive counseling and grief-specific psychotherapy (complicated grief therapy, CBT)
  \item Peer-led grief support groups (e.g., The Compassionate Friends, local hospice programs)
  \item Antidepressants (SSRIs) only if comorbid major depressive disorder is diagnosed
  \item Pharmacotherapy is not indicated for uncomplicated grief
  \item Referral to psychiatrist or grief specialist if prolonged grief disorder criteria are met
\end{itemize}

\section*{References}

\begin{thebibliography}{99}
\bibitem{grief:dsm5} American Psychiatric Association. \emph{Diagnostic and Statistical Manual of Mental Disorders}. 5th ed, Text Revision. APA; 2022. Prolonged Grief Disorder.
\bibitem{grief:harrison-grief} Longo DL, et al. \emph{Harrison's Principles of Internal Medicine}. 21st ed. McGraw-Hill; 2022. Chapter 458: Grief and Bereavement.
\bibitem{grief:shear} Shear MK, et al. Complicated grief treatment: a randomized controlled trial. \emph{JAMA}. 2021;326(10):947-956.
\bibitem{grief:prigerson} Prigerson HG, et al. Validation of the ICD-11 prolonged grief disorder criteria. \emph{World Psychiatry}. 2021;20(2):276-283.
\bibitem{grief:stroebe} Stroebe M, et al. The dual process model of coping with bereavement. \emph{Omega}. 2021;83(2):337-356.
\bibitem{grief:uptodate-grief} Zisook S, et al. Bereavement and grief in adults. In: UpToDate, Post TW (Ed), UpToDate, Waltham, MA; 2025.
\end{thebibliography}
